\inicializarSeccion{Explicación de campos mostrados}

\tab En las figuras \ref{fig:casoOchoCinco}, \ref{fig:casoOchoDiez} y \ref{fig:casoOchoVeinte}, podemos ver como parecen actuar como una carga puntual en el sentido de que, mientras mas lejos estas cargas estén, el campo electrico se vuelve mas débil. Sin embargo, desde estas figuras tambien podemos ver algo un poco mas interesante. Las cargas que estan en el centro de las figuras tienen casi un campo eléctrico nulo, esto debido a que las cargas positivas y negativas se anulan entre si; mientras mas alejadas las cargas del centro, el campo eléctrico en para la zona perpendicular a las lineas de carga se vuelve mas fuerte. 

\tab Analizando las figuras \ref{fig:casoTresCinco} y \ref{fig:casoVeinteCinco}, podemos ver que el campo elétrico; en el instante en la que tomamos el campo eléctrico, no necesariamente se vuelve mas fuerte a comparación de ambos casos. Esto se debe a que a que en ese preciso instante en donde el campo existe, las cargas todavia no se han separado lo suficiente como para iniciar una reacción en cadena para moverlas a una distancia mas lejana y por ende con mas fuerza.

\tab Algo muy interesante que se puede observar en las figuras \ref{fig:casoVeinteTresQuince} y \ref{fig:casoDosDiezQuince} es que se genera un campo eléctrico irregular con los huecos que tienen las lineas que contienen menos cargas. Este campo eléctrico irregular es la base por la cual nosotros podemos empezar a generar campos electricos candidatos a la dielectroforesis ya que al tenerlas en una misma diferencia entre una carga y otra, inevitablemente se generará un campo eléctrico en donde no todos los vectores se van a cerrar. Usando este principio, podemos empezar a mover biomoleculas que esten cargadas a través de un campo eléctrico.